\documentclass{report}
\usepackage[letterpaper,margin=35mm]{geometry}
\usepackage{parskip}
\usepackage{amsmath}
\usepackage{amssymb}

%re-format "\chapter{}" to my liking
\makeatletter
\def\@makechapterhead#1{{\vspace*{-2\baselineskip}
  \parindent \z@ \raggedright \normalfont
  \ifnum \c@secnumdepth >\m@ne {\huge\bfseries\thechapter.} \fi
  \interlinepenalty\@M \Huge \sffamily #1\par\nobreak \vskip 40\p@
}}
\makeatother

\pagestyle{plain}
\newcommand{\vbreak}{\vspace{4ex}}
\newcommand{\vextra}{\vspace{2ex}}
\newcommand{\QED}{\quad\blacksquare}
\newcommand{\Imax}{\ensuremath{I_\text{max}}}
\DeclareMathOperator{\atan}{atan}
\DeclareMathOperator{\err}{err}

\begin{document}
\chapter{Identities equivalent to $\atan(1)$}
A demonstration that
\begin{equation}\label{atan+10-239-515}
  8 \atan\tfrac{1}{10} -  \atan\tfrac{1}{239} - 4 \atan\tfrac{1}{515} =  \atan 1
\end{equation}

\vextra
First, recall the sum-of-angles identity for the tangent function:
\begin{equation}
  \tan(x+y) = \frac{\tan x + \tan y}{1 - (\tan x)(\tan y)}
\end{equation}
Furthermore, by the nature of inverse functions,
\begin{equation}
  x = \atan(\tan x) = \tan(\atan x)
\end{equation}
Therefore:
\begin{equation}
  \begin{aligned}
     \atan x + \atan y
       &= \atan( \tan( \atan x + \atan y ) ) \\
       &= \atan \frac{\tan(\atan x) + \tan(\atan y)}{1 - \tan(\atan x) \tan(\atan y)} \\
       &= \atan \frac{x + y}{1 - x y}
  \end{aligned}
\end{equation}
And, as a special case,
\begin{equation}
  \begin{aligned}
    2 \atan x &= \atan x + \atan x\\
              &= \atan \frac{2x}{1-x^2}
  \end{aligned}
\end{equation}

\vbreak
Thus:
\begin{equation}
  \begin{aligned}
    &8 \atan \tfrac{1}{10}  \\
       &\quad= 4 \atan \tfrac{2/10}{1 - 1/100}
	 \quad= 4 \atan \tfrac{20}{99} \\
       &\quad= 2 \atan \tfrac{40/99}{1 - 400/9801}
	 \quad= 2 \atan \tfrac{3960}{9401} \\
       &\quad= \atan \tfrac{7920/9401}{1 - 15681600/88378801}
	 \quad=   \atan \tfrac{74455920}{72697201}
  \end{aligned}
\end{equation}
And
\begin{equation}
  \begin{aligned}
    &4 \atan \tfrac{1}{515} \\
      &\quad= 2 \atan \tfrac{515}{132612} \\
      &\quad= \atan \tfrac{136590360}{17585677319}
  \end{aligned}
\end{equation}
So
  \begin{align*}
    &8 \atan\tfrac{1}{10} - \atan\tfrac{1}{239} - 4 \atan\tfrac{1}{515} \\
      &\quad= \atan\tfrac{74455920}{72697201} - (\atan\tfrac{1}{239} + \atan\tfrac{136590360}{17585677319}) \\
      &\quad= \atan\tfrac{74455920}{72697201} + \atan\tfrac{-1758719}{147153121} \\
      &\quad= \atan 1
      \QED
  \end{align*}

\vbreak
Another identity that could be used is:
\begin{equation}\label{atan+5-239}
   \atan 1 = 4 \atan\tfrac{1}{5} - \atan\tfrac{1}{239}
\end{equation}
This has the advantage of only requiring two passes,
but has the disadvantage that
the series for $\frac{1}{5}$ converges more slowly
than those for $\frac{1}{10}$ and $\frac{1}{515}$ combined.
The proof of correctness may be obtained in a manner very similar
to the one shown above,
and will not be spelled out here.

[Note that these Taylor-series based computations are not
state-of-the-art for computing absurd quantities of digits of tau or pi;
the trillion-plus digit record holders use a hypergeometric series
developed by the brothers David and Gregory Chudnovsky,
which cranks out about 15 digits per term computed for the series.
Another approach involves refinements of the Gauss-Legendre algorithm,
such as one by Richard Brent and Eugene Salamin (1976),
which \emph{doubles} the number of accurate digits with each iteration.
But these approaches are not as easy to understand or implement as
Taylor-series based approaches like the one used in this program.]

Putting this all together and running calculations
using multi-precision arithmetic with ``digit''s of base \texttt{BASE}
is how this code accomplishes its task.


\chapter{Error Analysis}
Let
\begin{equation}\label{function-f}
  f(x, n) = \frac{x^{2n+1}}{2n+1}
            \quad  [1 \le \tfrac{1}{x}; \text{ i.e., } 0 \le x \le 1]
\end{equation}
Then the Taylor-Maclaurin series for $\atan x$
can be written as:
\begin{equation}
  \atan x = \sum_{i=0}^{\infty} \,(-1)^i\, f(x, i)
\end{equation}

Recall that if $a_i$ is a monotonically non-increasing sequence
with $\lim_{i\to\infty} a_i = 0$,
then the sum $\sum_{i=0}^\infty (-1)^i\, a_i$ converges;
truncating this sum after term $n$
will have an error whose magnitude is at most $a_{n+1}$.
%((** double-check that all criteria are named and met **)
Thus for $\atan x$, provided that $1 \le \frac{1}{x}$,
the error incurred by truncating the series after the $n$-th term is:
\begin{equation}\label{partial-f-a-n}
  \err(n)
    \le f(x,n+1) = \frac{x^{2 n+3}}{2 n+3}
    <   \frac{x^{2 n+1}}{2 n+1} = f(x,n)
\end{equation}

We will assume that $a=1/x$ is the worst-case choice
for the relations being considered.
Substituting into \eqref{partial-f-a-n} gives
\begin{equation}\label{constraint-1}
  \err(n) < f(\tfrac{1}{a},n)
\end{equation}

\vbreak
We want our error term $\err(n)$ to be less than $B^{-D}$,
for some specified base $B$ and
to-be-determined base-$B$-digit-count $D$:
\begin{equation}\label{target}
  \err(n) < B^{-D}
\end{equation}

Furthermore, we want to fit intermediate results
into an integer which cannot exceed some maximum \Imax.
Each word of our multi-precision vector is
modulo-reduced by $2 n+1$ on outer-loop iteration $n$.
Each step computes
\[ r_i B + t_i \le \Imax \]
at word-index $i$,
for carry-down residue $r_i$ and previous-term $t_i$,
each of which is in the range $0 \le \star \le 2 n+1-1 = 2 n$,
thus the worst-case intermediate result is
\( 2\, n B + 2\, n \le \Imax \).
Solving for n yields:
\begin{equation}\label{n-terms}
  n \le \frac{\Imax}{2\,(B+1)}
\end{equation}

\vbreak
Our goal is to find the largest $D$ that simultaneously satisfies
relations \eqref{constraint-1}, \eqref{target}, and \eqref{n-terms}.
It is acceptable to find a value $D'$
that is less than the actual maximum,
if an approximation makes life easier,
but we must never allow the resultant $D'$ to be bigger than $D$
and we prefer not to underestimate $D$ by \emph{too} much.

\vextra
Now equations \eqref{constraint-1} and \eqref{target}
do not have a natural ordering,
but given our goal of wanting to derive a ``safe'' value of $D$,
we choose to combine them with this ordering:
\begin{equation}\label{target-inequality}
  \err(n) < f(\tfrac{1}{a},n) < B^{-D}
\end{equation}
For brevity, let
\begin{equation}\label{R-def}
  R = \frac{\Imax}{B+1}
\end{equation}
Combining substitutions of \eqref{n-terms} and \eqref{R-def}
into equation~\eqref{target-inequality} gives:
  \[ f(\tfrac{1}{a}, \tfrac{R}{2}) < B^{-D} \]
Solving for $D$ yields:
\begin{equation}
  D < \frac{ -\log f(\tfrac{1}{a}, \tfrac{R}{2}) }{ \log B }
\end{equation}
Using $(\tfrac{1}{a})^x = a^{-x}$,
$\log x/\log B = \log_B x$,
and the definition of $f$ given in equation~\eqref{function-f}:
\begin{equation}
	D
	< \frac{ -\log \frac{a^{-(2(R/2)+1)}}{2(R/2)+1} }{ \log B}
	= (R+1) \log_B a + \log_B (R+1)
\end{equation}
Converting the digit count to base-10 digits
involves multiplying by $\log_{10} B$.
Because
  \[ (\log_B x)(\log_{10} B)
     = \frac{\log_{10} x}{\log_{10} B} \log_{10} B
     = \log_{10} x
  \]
we can determine that the constraint on the number of decimal digits $d$ as
\begin{equation}\label{d-constraint}
  d = D \log_{10} B < (R+1) \log_{10} a + \log_{10}(R+1)
\end{equation}

\vbreak
To take a concrete example,
let us assume that our worst-case is $a=5$
(as when evaluating the relation in equation~\eqref{atan+5-239}),
that $B=10^{12}$,
and that $\Imax=2^{63}-1$;
then $R = \frac{2^{63}-1}{10^{12}+1}$
and we get
  \[ d
     < \left(\tfrac{2^{63}-1}{10^{12}+1} + 1\right) \log_{10} 5 + \log_{10} \left(\tfrac{2^{63}-1}{10^{12}+1} + 1\right)
     \approx 6\,446\,868
  \]
So for these parameters we want to ensure that no more than
$6\,446\,868$ decimal digits are allowed to be computed and printed,
in order to avoid the risk of errors from integer overflow
(whether they are silent as in C or panicking as in Rust).

\vbreak
Some further adjustments:
\begin{itemize}
  \item To re-iterate, we don't need to be precise in our upper bound;
    lowering by a few digits is acceptable.
  \item Hence it is okay to drop the $\log_{10}(R+1)$ term
    in equation~\eqref{d-constraint};
    even in the situation with
    $\Imax = 2^{64}$ and $B=10$,
    which will give a fairly large $R=\Imax/(B+1)$,
    we would lose fewer than 20 digits
    in our estimation of $d$.
  \item We are assuming that our worst case happens with $a=5$;
    $\log_{10}(5) > 339/485$
    with an error in the ballpark of one part per million.
  \item To minimize user confusion, we want to round the ``maximum''
    down to the nearest multiple of $\log_{10} B$,
    since we will always be outputting
    an integer number of base-$B$ ``digits''.
\end{itemize}
All of this can be accomplished with integer-only calculations
by using the following Rust function.
Here
\texttt{Xword} is the signed-integer type used for calculations,
\texttt{Xword::MAX} represents $\Imax$,
$\texttt{WORDDIGITS} = \log_{10} B$;
we require that \texttt{WORDDIGITS} be an integer,
hence that we end up with $\texttt{BASE} = B$ below:
\begin{verbatim}
     const BASE: Xword = (10 as Xword).pow(WORDDIGITS as u32);
     fn max_digits() -> Xword {
       let d = (Xword::MAX / (BASE+1) + 1) * 339 / 485;
       d - d % (WORDDIGITS as Xword)
     }
\end{verbatim}

\vextra
Note that if we use $a=10$
(such as when using equation~\eqref{atan+10-239-515})
instead of $a=5$
(as with equation~\eqref{atan+5-239}),
then instead of using $\log_{10}5 \approx 339/485$
we would use $\log_{10} 10 = 1$,
which both simplifies the equation and allows for the safe
computation of more digits.
On the other hand, an attempt to compute $\atan(1)$ directly
using the Taylor series would involve multiplying by
$\log_{10} 1 = 0$ --- thus we would
be limited to the fairly small
quantity of $\log_{10}(R+1)$
digits accurately computed
before encountering integer overflow problems
in some intermediate calculation;
this is one of the reasons we seek and use a mathematically
equivalent calculation which is more numerically stable
(the other reason is that even if we had infinite-precision variables,
the rate of convergence for a raw $\atan(1)$ is much too slow).

\end{document}
